\label{sec:experiments}
% >> TODO - add appendix with PDDL domains, generator params., transformer hyperparameters (num layers, batch size, lr, bce/goal-bce weight, focal coeff, dropout,), num predicates LS Transformer for each domain (GT conf.), <Explain how parameters \theta of LS Transformer preconds and effects are initialized! Also explain gumbel sigmoid temp> etc.
% Also --weight-init-method z (how LS Transformer params are initialized)
Our experiments address the following research questions:
\begin{itemize}
    \item[\textbf{RQ1:}] How do the LS and LB Transformers compare against standard transformer baselines?
    \item[\textbf{RQ2:}] How does symbolic alignment affect performance under final-state supervision, and how does the supervision setting affect the LS Transformer?
    \item[\textbf{RQ3:}] Can the trained transformers yield explicit \strips models that support planning?
    \item[\textbf{RQ4:}] Do the learned models extrapolate to longer traces and larger object sets?
\end{itemize}
%We explain our experimental setup below.
\paragraph{Baselines.}
We compare our models with two softmax-attention Transformer baselines:
\textbf{Sinusoidal}, using sinusoidal positional encodings
\citep{transformers}, and
\textbf{RoPE}, using rotary position embeddings
\citep{su2024roformer}, which can improve length generalization
\citep{spies2025transformers}.
Both use the same tokenization, randomized object embeddings, and network
dimensions as LB.
Positional information is applied to half of each token embedding to separate
it partially from the randomized object representation.
% \paragraph{Baselines.}
% We compare our proposed models against two standard baselines: \textbf{sinusoidal}, the original soft-attention transformer \citep{transformers}, with softmax aggregation and sinusoidal positional encodings; and \textbf{RoPE}, a soft-attention transformer with softmax aggregation and Rotary Position Embeddings \citep{su2024roformer}, which have been shown to improve generalization on longer sequences \citep{spies2025transformers}. 
%In both, positional information is applied only to half of each token embedding, to partially separate positional information from the randomized object embeddings.
\paragraph{Planning domains.}
We evaluate on standard automated-planning domains
\citep{long20033rd,fern2011first,taitler20242023}:
\textbf{\bw} (block rearrangement),
\textbf{\fr} (car transport by ferry),
\textbf{\lgt} (package delivery),
\textbf{\np} (sliding tiles), and
\textbf{\mz} (grid navigation).


% > Removed bullet points for space
\iffalse
In both, positional information is applied only to half of each token embedding, to partially separate positional information from the randomized object embeddings.

\begin{itemize}
    \item \textbf{Sinusoidal}: the original soft-attention transformer architecture proposed in \citep{transformers}, using softmax aggregation and sinusoidal positional encodings.
    \item \textbf{RoPE}: a soft-attention transformer with softmax aggregation and Rotary Position Embeddings \citep{su2024roformer}, which have been shown to improve generalization on longer sequences \citep{spies2025transformers}.
\end{itemize}

\paragraph{Planning domains.}
We perform experiments on five planning domains widely used in the AP community \citep{long20033rd,fern2011first,taitler20242023}:
\textbf{\bw}, where a set of blocks must be reassembled into a target configuration with a robotic gripper;
\textbf{\fr}, where a ferry is used to transport cars between different ports;
\textbf{\lgt}, a logistics task where a fleet of trucks and airplanes is used to deliver packages across locations in different cities;
\textbf{\np}, a puzzle where numbered tiles must be slid into their goal positions;
\textbf{\mz}, where an agent must traverse a grid-based maze to reach a goal location.
\fi

\paragraph{Data generation.}
For each domain, we generate 100 problems for each of the training, test, and planning splits using randomized generators.
Test and planning problems contain more objects than training problems. Training and test problems generate random traces whose numbers of domain actions are sampled uniformly from $[1,D]$. We generate $10^5$ training traces with $D=50$ and
$2\cdot10^5$ test traces with $D=200$. Exact object counts are reported in Appendix~\ref{app:problems}.

Under applicability supervision, each training action is sampled as applicable or inapplicable with equal probability. At test time, the prefix contains only applicable actions and only the final domain action may be inapplicable. Under final-state supervision, all domain actions are applicable, and each trace contains ten ground-atom queries sampled uniformly from all ground atoms over $O$.

\paragraph{Training and evaluation.}
The LS and LB Transformers are trained with Adam~\citep{liu2019variance}. We use $10^5$ training steps for LS and $5\cdot10^5$ for LB, increasing the latter to $7.5\cdot10^5$ in \lgt. Every $10^4$ steps, we compute accuracy on the full training set and retain the checkpoint with the highest training accuracy. We run each configuration with five random seeds and report mean performance, with the result of the seed achieving the highest training accuracy in parentheses.

Trace accuracy is the percentage of traces for which all supervised labels are predicted correctly. Under applicability supervision, these are the domain-action labels. Under final-state supervision, they are the final-state query labels. LS predictions use the binarized parameters $\bar{\theta}$, while LB
predictions are thresholded at $0.5$.
For RQ4, we vary trace length and object count, keeping the other dimension at its largest training value. We evaluate $10^4$ traces per setting and report the seed with the highest training accuracy.
We extract \strips models using the procedure associated with each supervision setting.
Planning accuracy is the percentage of planning problems for which the planner returns, within the search budget, a plan that is valid under the ground-truth model. All experiments run on an NVIDIA L40S GPU and an Intel Xeon Gold 6542Y CPU. Appendix~\ref{app:configuration} provides additional configuration details and model hyperparameters.

%\paragraph{Training and evaluation.}
%Both architectures are trained with RAdam \citep{liu2019variance}, using $10^5$ steps for the LS Transformer and $5\cdot10^5$ for the LB Transformer, increasing to $7.5\cdot10^5$ steps in \lgt for the latter.
%We perform validation every $10^4$ training steps, computing accuracy on the whole training dataset. After training, the checkpoint with the highest accuracy is tested. Accuracy is measured as the percentage of traces correctly classified, where a trace is considered correct only if the transformer correctly predicts the applicability label of each action in the trace. For the LS Transformer, predictions are obtained using the binarized parameters $\bar{\theta}$, whereas for the LB Transformer predictions are rounded to $\{0,1\}$.
%We extract \strips models from both architectures for planning, with the procedure corresponding to the supervision setting. Then, planning accuracy is computed as the percentage of planning problems correctly solved with the extracted model, i.e., those for which the planner outputs a valid plan.
%Finally, each experiment is conducted on a machine with an NVIDIA L40S GPU and an Intel Xeon Gold 6542Y CPU.
% If we don't say it before, say that we use GBFS FF and pymimir for planning!